\documentclass[letterpaper,12pt]{article}

\usepackage{
  amsmath,
  amsthm,
  amsfonts,
  natbib,
}

\bibliographystyle{plainnat}

\usepackage{hyperref}
\usepackage{cleveref}


\usepackage{todonotes}
%\usepackage[disable]{todonotes}

\title{Traceback ideas sketch}
\author{AA,LL,DR,AT,AY}

%%% macros
\DeclareMathOperator{\EX}{\mathbb{E}}
\DeclareMathOperator{\Hash}{\textsf{Hash}}
\DeclareMathOperator{\Enc}{\textsf{Enc}}
\DeclareMathOperator{\Dec}{\textsf{Dec}}
\DeclareMathOperator{\Sig}{\textsf{Sig}}
\DeclareMathOperator{\Ver}{\textsf{Ver}}
\DeclareMathOperator{\AddCount}{\textsf{AddCount}}
\DeclareMathOperator{\GetCount}{\textsf{GetCount}}
\newcommand{\pubkeyserv}{\textsf{pk}}
\newcommand{\seckeyserv}{\textsf{sk}}
\newcommand{\symkeyserv}{\textsf{key}}
\newcommand{\softO}{\tilde{O}}
%\newcommand{\softO}{O\tilde{~}}
\newcommand{\anote}[1]{{\sf (\textcolor{red}{Arkady's Note: {\sl{#1}}} {\sf EON)}}}


\begin{document}
\maketitle

Goal: Improve the ability of end-to-end encrypted messaging providers to
limit disinformation campaigns on their services, whereby a false or
manipulated story is shared via forwarding across the platform.

Specifically, we want to extend and improve on the ideas of
\citep{TMR-traceback-ccs19} by requiring a \emph{threshold} of $k$
distinct users to trigger a traceback (or other action).

Other possible goals are:
\begin{enumerate}
  \item Allow a range of \emph{actions} to include:
  \begin{enumerate}
    \item banning or suspending the account of the content originator
    \item sending a notification to all who have already received the
    content
    \item showing a notification/warning to all who view the content
    later (regardless of when they received it)
    \item preventing any additional forwards of the content
  \end{enumerate}
  \item Disclose less information than a full tree/path traceback, for
  example just the content of the message, or just the originator.
  \item Support one or more independent auditors who will review reports
  and decide on the appropriate action to take (if any)
  \item Maintain privacy, the the extent possible, in the face of
  collusion and possibly malicious actions by some group of users, the
  service provider, maybe even the auditor
\end{enumerate}

\subsection{Notation and Nomenclature}

\begin{description}
  \item[Service provider (SP)] The web service communicating with users and
    handling messages between them, e.g., WhatsApp or Signal.
  \item[Auditor] An trusted party who reviews verified reports and
    decides what action needs to be taken (if any).
  \item[User] An ordinary user of the messaging system, utilizing end-to-end
    encryption and possibly receiving or sharing misinformation.
  \item[Message] A piece of data (text, image, whatever) sent from a
    single user to one or more other users in encrypted form.
    Should ordinarily be indistinguishable
    from any other message of roughly the same length to anyone that is
    not the sender or the recipient.
  \item[Forward]
  \item[Original message]
  \item[Originator] The user who first sent an original message
  \item[UserId] Each user has an immutable, unique, public identifier
    number, basically a public key whose corresponding private key can
    be used for signing or other more exotic things. The list of users
    should be public, or at least checkable by the auditor, so that the
    service provider can't too easily create new fake user accounts.
  \item[MessageId] An immutable value known only to those who have
    received a message (or a forward of it).
    Something like a deterministic hash of (some combination of) the
    message content, originator, and first send time, plus maybe some
    salt.
  \item[Report] The action a user takes to complain that a message
    is/may be misinformation. In issuing a report, the user is willing
    to waive the privacy of the message contents as well as (perhaps)
    their identity, at least to the auditor.
  \item[Threshold] The minimum number of users, written as $k$, who must
    independently report the same message before it is sent to the
    auditor to determine what action to take.
\end{description}

\subsection{Security goals}

There are lots of competing privacy and security goals.

Privacy:
\begin{itemize}
  \item For any message which is not reported,
    and is not sent to any colluding user,
    no one else (SP, other users, auditor) should be able to see the
    contents of a message, or even if it is a forward or not.
  \item The same should hold for any message which is reported by fewer
    than $k$ users (the threshold)
  \item The SP should not learn anything about which message a user
    reports, until the threshold is reached.
  \item The auditor \emph{should} learn the message contents and the
    identities of the $k$ reporters, but only after the threshold has been
    reached.
\end{itemize}

Security:
\begin{itemize}
  \item The auditor must approve any punitive action for disinformation.
    That is, the SP should convince all users that experience a punitive
    action that the auditor really approved it.
  \item Any collusion between some users and/or SP can only trick the
    auditor into approving an action on a
    message which has been sent to one of those colluding users, and only
    if the number of colluding users is at least $k$.
  \item No user can disguise or alter the originator of a message being
    forwarded.
\end{itemize}

Also regarding performance expectations, the users and auditor should
have minimal long-term storage and computation requirements.

\subsection{Collaborative counting Bloom filter (CCBF)}\label{sec:ctbf}

Here we describe the CCBF data structure which we will use to track
complaints. The CCBF is somewhat similar to a count-min sketch in functionality:
it allows incrementing the count of any item, and performing a
\emph{point query} to estimate the current count of any item.
Unlike count-min sketches, the CCBF does not support any other
type of query and exhibits a relatively wide two-sided error bias on the
count estimates returned.

The CCBF supports two operations:
\begin{itemize}
  \item $\AddCount(i,x)$: Adds identifier $i$ to the set of witnesses for item $x$
  \item $\GetCount(x)$: Computes an estimate on the number of distinct
    witnesses $i$ in the set of counts for item $x$
\end{itemize}

Of course, a CCBF could be instantiated by simply storing a list of
$(i,x)$ pairs, but we want a construction which --- like all Bloom
filter variants --- trades improved memory efficiency for decreased
accuracy in the $\GetCount$ operation.

The construction is relatively simple at a high level, parameterized by
three values $s,u,v$ which must be fixed at instantiation time:
\begin{itemize}
  \item The storage is a length-$s$ table of bits $T$, initially all zeros.
  \item Each identifier $i$ is associated with a random size-$u$ subset
    of table indexes $\{0,1,\ldots,s-1\}$, denoted $U(i)$.
  \item Each item $x$ is associated with a random size-$v$ subset of
    table indices, denoted $V(x)$.
  \item The $\AddCount(i,x)$ operation computes the intersection of the
    identifier and item index sets, $U(i)\bigcap V(x)$, and checks these
    locations in the table. If possible, exactly one of these locations
    in the table is flipped from 0 to 1.
  \item The $\GetCount(x)$ operation looks up all table locations in
    $V(x)$ and counts how many of the $v$ locations are set to 1. This
    is fed into a calculation, based on the total number of set bits in
    the entire table, to produce an estimate of the current count.
\end{itemize}

\subsubsection{Elementary approximations}

We begin with some standard building-block approximations.

First, we can approximate the function $\exp(x)$ around the origin
with straight lines.

The following inequalities hold for any constant $\alpha\ge 0$ and any
real $x$ with $0\le x \le \alpha$:
\begin{eqnarray}
  \exp\left(\frac{\ln(1+\alpha)}{\alpha} x\right)
  \le 1 + x
  \le \exp(x)
  \le 1 + \left(\frac{\exp(\alpha)-1}{\alpha}\right) x
  \label{eqn:inequal-pos}
\\[15 pt]
  \exp\left(-\frac{\ln\frac{1}{1-\alpha}}{\alpha} x\right)
  \le 1 - x
  \le \exp(-x)
  \le 1 - \frac{1 - \exp(-\alpha)}{\alpha} x
  \label{eqn:inequal-neg}
\end{eqnarray}

We can invert the inequalities on the ends of the lines above to
obtain:
\begin{eqnarray}
%    \exp\left(\tfrac{x}{2}\right) \le 1+x
%      & 0 \le x \le 2.5128 \\
  \exp(x) \le 1 + 2x,\quad
  0 \le x \le 1.2564
  \label{eqn:bound2-exp-pos}
\\[10 pt]
  \exp(-2x) \le 1-x,\quad
  0 \le x \le 0.7968
  \label{eqn:bound2-exp-neg}
%   \exp(-x) \le 1 - \tfrac{x}{2},\quad
%   0 \le x \le 1.5936
\end{eqnarray}

We also need approximations for (partial) harmonic series of the form
\(\frac{1}{a} + \frac{1}{a+1} + \cdots + \frac{1}{b}\). Basic left and
right Riemann sums give, for two integers \(a \le b\),
\[\sum_{i=a+1}^{b} \frac{1}{i} \le \ln b - \ln a \le
\sum_{i=a}^{b-1}\frac{1}{i},\]\anote{Why do the indices differ in the sums on the left and right?}
which is inverted on both sides to obtain
\begin{equation}\label{eq:harm-bound}
\ln\frac{b+1}{a} \le \sum_{i=a}^b \frac{1}{i} \le \ln\frac{b}{a-1}.
\end{equation}

\subsubsection{Calculation of $\GetCount$}

Recall that $s$ is the total table size, $u$ is the number of slots per
identifier, and $v$ is the number of slots per item.

\paragraph{Probability to fill in table during $\AddCount$.}

Suppose $\AddCount(i,x)$ is called at a time when
$v_0$ out of the $V(x)$ indexed bits are set to
1. We want to know the probability that this call to $\AddCount$ sets
one more bit out of $V(x)$.

Consider the intersection of the $(v-v_0)$ unset indices of $V(x)$, with
all $u$ identifier indices $U(i)$. Assuming $\AddCount(i,x)$ has not
been called already, we can consider these as two independent subsets of
an $s$-element set. The exact probability that they do \emph{not}
intersect is
\[\frac{\binom{s-v+v_0}{u}}{\binom{s}{u}}
  = \frac{(s-v+v_0)! (s-u)!}{s! (s-u-v+v_0)!}
  = \prod_{j=0}^{v-v_0+1} \frac{s-u-j}{s-j}.\]
Note that this probability can be computed exactly in $\softO(s)$ time.

An unconditional upper bound is obtained from \eqref{eqn:inequal-neg}:
\[ \prod_{j=0}^{v-v_0+1} \frac{s-u-j}{s-j}
  \le \left(1 - \frac{u}{s}\right)^{v-v_0}
  \le \exp\left(-\frac{u(v-v_0)}{s}\right)
\]

% XXX: This is the derivation, replaced with explicit ratios below.
% Let $c_1\ge v/s$ and $c_2 \ge u/s$. Then, provided $c_1,c_2$ are
% sufficiently small,
% we can apply \eqref{eqn:inequal-neg} to obtain a lower bound as
% well:
% \begin{eqnarray*}
%   \prod_{j=0}^{v-v_0+1} \frac{s-u-j}{s-j}
%     &\ge& \left(1 - \frac{u}{s-v+v_0+1}\right)^{v-v_0} \\
%     &>& \left(1 - \frac{u}{s-v}\right)^{v-v_0} \\
%     &\ge& \left(1 - \frac{u}{(1-c_1)s}\right)^{v-v_0} \\
%     &\ge& \exp\left(- \frac{\ln\left(1 + \frac{c2}{1-c_1-c_2}\right)}{c_2}
%       \frac{u(v-v_0)}{s}\right)
% \end{eqnarray*}

Provided that $u \le 0.5s$ and $v \le 0.2s$,
we can apply \eqref{eqn:inequal-neg} to obtain a lower bound as
well:
\begin{eqnarray*}
  \prod_{j=0}^{v-v_0+1} \frac{s-u-j}{s-j}
    &\ge& \left(1 - \frac{u}{s-v+v_0+1}\right)^{v-v_0} \\
    &>& \left(1 - \frac{u}{s-v}\right)^{v-v_0} \\
    &\ge& \left(1 - \frac{u}{0.8 s}\right)^{v-v_0} \\
    &\ge& \exp\left(-2 \frac{u(v-v_0)}{s}\right)
\end{eqnarray*}

In summary, whenever $u\le 0.5s$ and $v \le 0.2s$, we have
\begin{equation}\label{eqn:prob-no-inter}
  \exp(-2u(v-v_0)/s) < \Pr[\text{no intersection}] < \exp(-u(v-v_0)/s).
\end{equation}

\paragraph{Expected time to fill $k$ slots for item $x$.}

Following from the previous estimate,
the expected number of calls to $\AddCount(\cdot,x)$ before the next
entry of $V(x)$ is actually filled follows a geometric distribution,
with expected value $1/(1 - \Pr[\text{no intersection}])$. Again, this
expected value can be computed exactly.

Using the bounds above, under the more restrictive conditions that
$8 \le u \le 0.5s$ and $uv \le 1.59s$, we can approximate this expected
value as
\begin{equation}\label{eq:exptries}
  0.5\frac{s}{u(v-v_0)} < \EX[\text{tries until next set bit}] <
  2\frac{s}{u(v-v_0)}.
\end{equation}

Now we want to ask, what is the total \emph{expected} number of calls to
$\AddCount(\cdot,x)$ that occur before some $k$ out of the $v$ total
indices in $V(x)$ are set to 1? By linearity of expectation, we can
simply sum up the expected number of tries to set each individual bit,
starting with the amount of bits set by \emph{other} calls to
$\AddCount(\cdot,y)$ with $y\ne x$.

Considering only the previous calls to $\AddCount(\cdot,y)$ with $y\ne
x$, suppose the resulting table has $s_0$ total bits set to 1. Then,
because the index set $V(x)$ is chosen independently at random, the
expected number of bits within $V(x)$ set by other calls to $\AddCount$
is exactly $vs_0/s$.

The total expected number of calls to $\AddCount(\cdot,x)$ needed to
have $k$ total bits set is therefore
\begin{equation}\label{eq:totalexp}
  \sum_{i=vs_0/s}^{k-1} \EX[\text{tries to set the $i+1$'st bit}],
\end{equation}
which again can be computed exactly.

We can approximate \eqref{eq:totalexp} using \eqref{eq:exptries}. First,
for an upper bound, the exact value of \eqref{eq:totalexp} is at most
\[\sum_{i=vs_0/s}^{k-1} \frac{2s}{u(v-i)}
  = \frac{2s}{u}\sum_{j=v-k+1}^{v-vs_0/s} \frac{1}{j}.\]

Using the upper bound in \eqref{eq:harm-bound} and the fact that
$s_0\ge 0$, we can approximate this
summation as

\[\sum_{j=v-k+1}^{v-vs_0/s}\frac{1}{j}
  \le \ln \frac{v - vs_0/s}{v-k} \le \ln\left(1 +
  \frac{k}{v-k}\right)\]

Finally, since \(\ln(1+x) \le x\) for any $x\ge 0$, incorporating this
into the previous inequality, the expected number of calls to
$\AddCount(\cdot,x)$ to have $k$ total bits set is at most
\begin{equation}
  \frac{2s}{u\left(\tfrac{v}{k}-1\right)}
\end{equation}

For a lower bound, we need to assume that the number of bits set $k/v$
is significantly larger than the ``background noise'' corresponding to
$s_0/s$. Specifically, assuming $\frac{k}{v}\ge 2\frac{s_0}{s}$, and
using the same inequalities as before we can derive a lower bound on
\eqref{eq:totalexp} of

\begin{eqnarray*}
\frac{s}{2u}\sum_{j=v-k+1}^{v-vs_0/s} \frac{1}{j}
  &\ge& \frac{s}{2u}\ln\frac{v-vs_0/s+1}{v-k+1} \\
  &\ge& \frac{s}{2u}\ln\frac{v+1-k/2}{v+1-k} \\
  &\ge& \frac{s}{2u}\ln\left(1 + \frac{k/2}{v+1-k}\right)
\end{eqnarray*}

%TODO check w.r.t. previous bounds
Now assuming \(0 \le \frac{k/2}{v+1-k} \le 2.5128\) and
\(v+1-k \ge 5\),
we can apply \eqref{eqn:bound2-exp-pos}
to further simplify the lower bound on the number of valid complaints
needed to fill in $k$ slots for the message:

\begin{eqnarray*}
\frac{s}{2u}\ln\left(1 + \frac{k/2}{v+1-k}\right)
  &\ge& \frac{s}{2u} \cdot \frac{k/4}{v+1-k} \\
  &\ge& \frac{s}{2u} \cdot \frac{k/5}{v-k} \\
  &=& \frac{s}{10u(\tfrac{v}{k}-1)}
\end{eqnarray*}

Putting this together with the upper bound, we have
\begin{equation}\label{eq:exp-fill}
  0.1 \frac{s}{u(\tfrac{v}{k}-1)}
  \le \EX[\text{valid complaints to fill $k$ bits}]
  \le 2 \frac{s}{u(\tfrac{v}{k}-1)}
\end{equation}

\subsubsection{Tail bounds}

Let $\lambda$ be a (statistical) security parameter. We need tail bounds
relative to $\exp(-\lambda)$ on the probability that:
\begin{itemize}
  \item Many valid complaints do \emph{not} trigger an audit; or
  \item Few valid complaints \emph{do} trigger any audit.
\end{itemize}

For both of these, we can consider a sequence of $c$ attempted
complaints by different users on the same message. By the principle of
deferred decisions, we can consider the probability each attempted complaint does
not succeed in writing to one of the message's $v$ slots, to be the
probability that a random size-$u$ subset
does not intersect with a given size-$(v-v_0)$ subset of all $s$ positions.
Recall from \eqref{eqn:prob-no-inter} that this is given by
\[
  \exp(-2u(v-v_0)/s) < \Pr[\text{no intersection}] < \exp(-u(v-v_0)/s).
\]

Now define indicator random variables $X_i$ to be $1$ if the $i$th
attempted complaint does not succeed, and $0$ otherwise. We are
interested in the total number $\bar{X}$ of failures among $c$
complaints:
\[
  \bar{X} = X_1 + X_2 + \cdots + X_c
\]

Each $X_i$ is in fact a Bernoulli distribution with probability $p_i$
bounded by \eqref{eqn:prob-no-inter}. These random variables are
\emph{not} independent, as the parameter $v_0$ is changing
depending on prior events when some message slots get filled in, and
this affects the probability $p_i$.

However, we can bound the probabilities $p_i$ and therefore the sum
$\bar{X}$ by the sum of i.i.d. Bernoulli random variables, and then use
Hoeffding's inequality to derive the required tail bounds.

First for the upper bound, we want to show that, for sufficiently many
attempted complaints $c$, the probability that no more than $k-1$ of
them succeed is very small. That is, we want to show that
\[\Pr[\bar{X} \ge c-k+1] \le \exp(-\lambda).\]

For this, we start with the following assumptions:
\begin{enumerate}
  \item \label{ass:kv}
    \(2\frac{s_0}{s} \le \frac{k}{v} \le 0.5\)
  \item \label{ass:k}
    \(k \ge 5\)
  \item \label{ass:u}
    \(8 \le u \le 0.5s\)
  \item \label{ass:uv}
    \(\frac{s}{uv} \ge 0.629\)
\end{enumerate}

In that case that fewer than $k$ of the $v$ message slots are filled
after the entire sequence of $c$ complaints, the probability of failure
for each attempted complaint is maximized when $k-1$ slots have been
filled. From \eqref{eqn:prob-no-inter}, this gives an upper bound on the
probability of each failure of
\[p_i \le \exp({-u(v-k+1)/s}).\]

Call this bound $q$, and define a new sequence of i.i.d. Bernoulli random variables
$Y_1,\ldots,Y_c$, which equal 1 with probability $q$ and 0 otherwise,
and $\bar{Y} = \sum_{i=1}^c Y_i$ as their sum. Because the probability
$p_i$ of each Bernoulli r.v. $X_i$ is bounded above by $q$, we conclude
that
\[\Pr[\bar{X} \ge c-k+1] \le \Pr[\bar{Y} \ge c-k+1]\]
and turn our attention to the latter probability.

Now let
\begin{equation}\label{eqn:c-upper-tail}
  c \ge 9\lambda\left(\frac{s}{uv}\right)^2 + 6\frac{ks}{uv}.
\end{equation}

We first claim that $3\frac{s}{uv} \ge \frac{1}{1-q}$, where
$q=\exp(-u(v-k+1)/s)$ is the upper bound on each $p_i$ as derived above.
By Assumption~\ref{ass:kv}, we have $v-k+1 > v-k \ge 0.5 v$, and
therefore $q \le \exp(-0.5 \frac{uv}{s})$. Because of
Assumption~\ref{ass:uv}, we know that $0.5\frac{uv}{s} < 0.795$, which
means we can apply the bound of \eqref{eqn:inequal-neg}
(with $\alpha\approx 0.5484$) to derive
\[q \le \exp(-0.5 \frac{uv}{s}) < 1 - \frac{1}{3} \frac{uv}{s},\]
and therefore $\frac{1}{1-q} < 3\frac{s}{uv}$ as claimed.

Applying this to \eqref{eqn:c-upper-tail}, we have
\begin{eqnarray*}
  c &>& \frac{\lambda}{(1-q)^2} + \frac{2(k-1)}{1-q} \\
  &=& 2\left(\frac{\sqrt{\lambda/2}}{1-q}\right)^2
    + 2\left(\frac{\sqrt{k-1}}{\sqrt{1-q}}\right)^2.
\end{eqnarray*}

Next, because \((a+b)^2 \le 2a^2 + 2b^2\) for any real numbers $a,b$, we
know that
\[\sqrt{c} > \frac{\sqrt{\lambda/2} + \sqrt{(k-1)(1-q)}}{1-q},\]
and therefore
\begin{eqnarray*}
(1-q)c
  &>& \left(\sqrt{\lambda/2} + \sqrt{(k-1)(1-q)}\right)\sqrt{c} \\
  &>& \sqrt{\lambda c/2} + k - 1.
\end{eqnarray*}

Rearranging gives $c-k+1 \ge \left(q + \sqrt{\lambda/(2c)}\right)c$,
which means we can finally apply Hoeffding's inequality over the sum of
i.i.d. Bernoulli r.v.'s $\bar{Y}$ to conclude that
\[\Pr[\bar{Y} \ge c-k+1] \le \exp(-\lambda).\]

From the discussion above, this proves that, except with negligible
probability, any $c$ attempted complaints on the same message will
always trigger an audit, with $c$ as in \eqref{eqn:c-upper-tail}.


\subsubsection{Oblivious CCBF}

A \emph{write-oblivious} data structure is one for which an adversary
who can see all updates to underlying memory storage cannot learn
anything more than the number of data structure (write) operations which
were issued (in our context the number of times $\AddCount$ was called).

The CCBF described above is already \emph{nearly} write-oblivious: In
the generic case, a call to $\AddCount$ results in exactly one bit in
the table being set to 1. No other $\AddCount$ operation will touch this
same bit, and it is indistinguishable from random. However, if the
index set intersection $U(i)\bigcap V(i)$ is empty, or if all indices in
this set are already 1, then no write occurs. This difference violates
write-obliviousness.

To avoid it, we need a cap $c$ on the total number of times $\AddCount$
will be called, a limit $q$ on how many times $\AddCount(i,\cdot)$
will be called for any identifier $i$, and a security parameter
$\lambda$.

Then the $\AddCount$ algorithm checks not only the table indices in
$U(i)\bigcap V(x)$, but also up to $(\lambda+q-1)$ additional indices from
$U(i)$. If the bit for any index in the actual intersection is 0, it is
set to 1 as usual. But otherwise, one of the bits from the additional
indices is changed from 0 to 1.

If this algorithm always succeeds, then it is write-oblivious since
every $\AddCount$ operation sets exactly 1 bit, which is never again
written to, and the index of this bit is indistinguishable from random.

However, catastrophic failure occurs if all of the bits at selected
indices are all 1. To ensure this happens with negligible probability,
we need only ensure the table size $s$ is at least $2c$, and the
identifier index set size $u$ is at least $(\lambda + q-1)$. In this case,
during any call to $\AddCount(i,\cdot)$, the set $U(i)$ will have at
most $q-1$ bits set from previous calls to $\AddCount(i,\cdot)$, and the
chance that any additional bit is 1 is at most $c/s$. Therefore the
chance that all $(\lambda+q-1)$ selected ``extra'' indices from $U(i)$
are 1, in any call to $\AddCount$, is at most $2^{-\lambda}$.

The CCBF can be further made \emph{completely} (read- and write-)
oblivious by replacing the table lookups with private information
retrieval (PIR). Note that the size of the table lookups (number of bits
checked) is at most $v+q+\lambda-1$ in any call to $\AddCount$ or
$\GetCount$, so the efficiency of the PIR operation depends only on
these parameters and the table size $s$.

\subsection{Solution based on CTBF}

Here we describe a solution to the message-reporting problem with
desired threshold $t$, using the oblivious CTBF data structure described
in the previous section.

\begin{itemize}
  \item The server stores a CTBF table (i.e., a length-$s$ bit array).

  \item Each user has an identifier $i$ and is assigned a random subset $U(i)$
    of table indices. This subset $U(i)$ should be known only to the
    user and the server. It could be chosen randomly by the server and
    sent to the user, or computed from some PRG from a key shared with
    each user and the server.

  \item On origination, each message $x$ is assigned a random size-$v$
    $V(x)$. In this case, the choice of subset must be random (even if
    the user is malicious), but the server cannot learn the subset. This
    can be achieved by having the originator commit to a random value
    $r_u$, the server write a second random value $r_s$ while signing
    the message hash, and then use their xor $r_u\oplus r_s$ as seed for
    a PRG that determines the subset. See \cref{sec:unforgeable} for
    details on how to construct the message.

  \item Forwarding a message just means sending the (signed, originated)
    message through the E2E messaging channel. If this channel is
    visible to the server, the forwarding user can first originate a
    dummy message in order to not reveal whether the nature of the
    message.

  \item Any honest user receiving a message should first check that it
    has a valid construction and that the signature verifies with the
    server's public key. If not, the message is discarded (and optionally,
    the server may be notified).

  \item To issue a complaint on message $x$, user $i$ first performs
    the $\AddCount(i,x)$ operation obliviously, using PIR for the
    lookups. Note, this is done authenticated with the server, so the
    server can check that user $i$ has not yet reached their quota $q$
    of complaints.

    The user then computes $\GetCount(x)$ (having already done
    the necessary lookups), and if the value exceeds the desired
    threshold $t$, the user reveals the plaintext message value, along
    with signed commitment to $V(x)$, to the auditor.

    The auditor (who may be the same as the server)
    can then examine the message contents, do fact-checking etc., and
    flag the message, ban the originator, etc.
\end{itemize}

Security and efficiency come from the signature/unforgeability
properties as outlined in \cref{sec:unforgeable}, and the CTBF
construction from \cref{sec:ctbf}.

\subsubsection{Parameters}

We assume this process is conducted over some fixed \emph{time window},
say 1 day. For that day, the system designer must choose values for:
\begin{itemize}
  \item $N$: Upper bound on total number of distinct users
  \item $m$: Upper bound on number of messages sent in the time window
  \item $c$: Upper bound on total number of attempted complaints within
    the time window
  \item $t$: \emph{Desired} threshold; this will be the expected number
    of honest complaints to trigger an audit.
  \item $q$: Quota on maximum number of writes (i.e., complaints)
    allowed per user within the time window
  \item $\lambda$: Security parameter
\end{itemize}

Based on these values and some performance and security criteria, the
scheme will determine values for:
\begin{itemize}
  \item $s$: Size of the table, in bits
  \item $u$: Number of slots each user is allowed to write
  \item $v$: Number of slots for each message
\end{itemize}

The server keeps track of a value $s_0\le c$ for how many table bits are
currently set to 1, and this value is shared with all users before they
issue a complaint; that is, $s_0$ is public information.

\subsection{Building block ideas}

\subsubsection{Probabilistic threshold}\label{sec:prob-thresh}

Each \emph{attempted} report may succeed or fail, depending on a
deterministic function of the MessageId and UserId. This deterministic
function should have a known probability distribution, so that each user
has a fixed probability $p$ of \emph{successfully} reporting a given
message.

For example, to achieve a probability $\tfrac{1}{4}$, the function could
be checking that the last two bits of the MessageId and UserId are the
same.

But ideally the function should know be known to the users or they could
maybe do some brute-force checks themselves. So either the MessageId
should be collaboratively chosen at message creation time, or the
reporting success function will be computed with some input from the
service provider.

The advantage is that, by setting $p=1/k$, then the \emph{expected}
number of attempted reports before a single successful report is $k$.
Unfortunately the \emph{variance} is also quite high; see Wikipedia page
on the geometric distribution.

\subsubsection{PIR and Bloom Filters}\label{sec:pir-bloom}

The SP maintains a Bloom Filter. (It would have to be cleared and
refreshed periodically to avoid too many false positives.)

Each MessageId deterministically maps to a set of $k-1$ Bloom filter
locations; these are where the signed (successful) reports will be
stored.

To issue a report:
\begin{enumerate}
  \item User computes the list of all $k-1$ locations for that
  MessageId.
  \item User prepares their own report, which is their UserId signature
  over the MessageId, encrypted with the Auditor's public key.
  \item User queries all $k-1$ locations via PIR with SP. This way, the
  service provider does not learn anything about which message is being
  looked up. Furthermore, a user only knows the locations for messages
  they have actually received, so can't learn extra information from
  ther results.
  \item If all $k-1$ locations are full, then the list of locations,
  plus this user's own report, are sent to the SP. The SP sends all
  these to the auditor to check the report and determine any action.
  \item If any of the $k-1$ locations is empty, the user inerts their
  own report at that location. Note that this does not reveal anything
  to SP, because that location has not been used previously.
\end{enumerate}

Note, this can be combined with the probabilistic threshold technique
(\cref{sec:prob-thresh}), to achive a smaller variance.

\textbf{Problems}: Handling overlap between the location sets in the
Bloom filter. The probability of overlap can be made small, but this
will also make the total size larger. Alternatively, somehow allow the
user to check whether each location holds a report \emph{for the same
MessageId}, and the SP holds a linked list of reports on each location;
but this might make the PIR too expensive or cause some other leakage.

\subsubsection{Unforgeable message payload}\label{sec:unforgeable}

Goals:
\begin{itemize}
  \item
    Every message has an author, the \emph{originator} who sent this
    message the first time.
  \item
    Every message sent (original or forward) should contain some
    ciphertext that stores the originator's identity. Either the server
    or the (honest) receiving client should be able to detect and
    discard any improperly-structured message that doesn't contain this.
  \item
    The originator \emph{must} include her own identity with a message
    she writes, and no one else can attach someone else's identity to
    a different message.
\end{itemize}

We can achieve this by using a normal cryptographic hash $\Hash$, a
symmetric encryption scheme $(\Enc,\Dec)$, and a signature scheme
$(\Sig,\Ver)$.

Each message has some plaintext
contents $m$, which will include the actual message that a client sees,
plus any other information (maybe stuff related to secret sharing, or
something else) that will be revealed to every client who receives the
message, but not by the server.

The server has an asymmetric key pair $(\pubkeyserv,\seckeyserv)$
for signing, and every client knows $\pubkeyserv$.
The server also has a secret symmetric key $\symkeyserv$.

(Note, the clients likely also have asymmetric key pairs for sending
messages, but we assume this is part of the existing E2EE messaging
system that we are building on top of.)

The protocol to create a new message
is as follows:

\begin{enumerate}
  \item Client $i$ chooses random salt $s$, computes a salted hash $h \gets \Hash(s,m)$,
    and sends $h$ to the server over an authenticated channel.
  \item Server computes a (randomized) encryption of client $i$'s
    identity, $e \gets \Enc_\symkeyserv(i)$, a signature
    $g \gets \Sig(h, e)$,
    and sends the tuple $(e, g)$ to the client.
  \item Client $i$ computes the tuple $c \gets (s,m,e,g)$, which will be
    the contents of the E2EE message actually sent to some other client
    $j$.
\end{enumerate}

Upon receiving an E2EE message $(s,m,e,g)$, client $j$ must check
$\Ver_\pubkeyserv(g,(\Hash(s,m),e))$ to ensure that this is an
unmodified message attached to the originator's identity.

A valid complaint (above the threshold\ldots) should reveal $(s,m,e)$ to
the server, from which the server can compute $\Dec_\symkeyserv(e)$ and
discover the identity of the message's originator. Note, the encryption
$e$ could also contain some other information that the server wants to
learn on a traceback, such as a timestamp.

To flag a message, the server can broadcast the hash $h=\Hash(s,m)$ to
all clients; any honest client who later receives this message will know
to display it with a warning.

\subsection{Discared ideas?}

\subsubsection{Multiple ORAMs}

Each user has an ORAM, all stored by the service provider. Each ORAM
contains a list of signed MessageId's, corresponding to previous
successful reports.

To check a message, a user looks through their ORAM for all messages
matching the current MessageId.

To issue a report, the user adds their signature on the MessageId to
every other ORAM (with help from Service Provider). Note that adding can
be done with a simple ``append'' operation on the ORAM stash.

When issuing, a user will see if the number of reports has reached the
threshold. In this case, the user sends all signatures to the Service
Provider and/or Auditor to take action.

Note, this can be combined with the probabilistic threshold technique
(\cref{sec:prob-thresh}), to achive a smaller variance.

\textbf{Problems}: Lots of storage for the SP (\# of users times \# of
reports). How to check against fake reports (maybe need some more
crypto).

% !TEX root = ideas.tex
%\newcommand{\anote}[1]{\todo[inline, size=\small, author=Arkady]{#1}}
\newcommand{\tuple}[2]{\langle #1, #2 \rangle }
\newcommand{\tup}[1]{\langle #1, #1 \rangle}
\subsubsection{Server-Checked Complaints}
\paragraph{A fatal flaw?}
In the process of writing this section, I realized that my approach has a fatal flaw.  The problem is with reusing the same check poly, $cp$, for all messages.  Now, when a party $u_i$ receives any two messages, he can compute two ciphertexts with the same value of $cp(i)$.  He can then subtract these homomorphically, to get a ciphertext with a 0 in the $cp$ slot.  This enables him to change the $p(i)$ value in any subsequent complaint without violating correctness of the check poly.  

\paragraph{Original Scheme:}
One problem that we have is clients submitting invalid complaints.  Ideally, we would like to allow the Server to check whether a complaint is valid without learning what message the complaint corresponds to.

Below is a proposed solution for this, though many details in particular, the originating procedure needs to be optimized.

\paragraph{Crypto Tools:}
We use a secret-key homomorphic encryption (HE) scheme that operates on tuples of values $(x_1,x_2)$.  The HE scheme needs to be able to compute additive homomorphism and one level of multiplications.
We need the following security properties (satisfied by many HE schemes)
\begin{itemize}
\item It is hard for someone not knowing $sk$ to produce a valid ciphertext from scratch
\item When given Enc(f(x)) produced by homomorphism, it is not possible to tell what $f$ was.  A stronger form of this, known as circuit privacy, requires that one (even with $sk$) cannot tell the difference between a freshly encrypted $ct$ and one produced via homomorphism.
\end{itemize}
\paragraph{Setup:}
\begin{enumerate}
\item $S$ generates $sk$ for HE.
\item $S$ picks a random degree $t$ check poly, $cp(x)$.  Let $cp_0,\ldots,cp_t$ be the coefficients of $cp$.
\item When a user $u_i$ joins, $S$ computes $Enc_{sk}(\tup{i} ,Enc_{sk}(\tup{i^2}),\ldots,Enc_{sk}(\tup{i^t})$ and gives these to $u_i$.
\end{enumerate}

\paragraph{Originating a message: $S$ and $u_j$ where $j$ is the originator's ID}
\begin{enumerate}
\item $S$ picks a random degree $t$ poly $p_S(x)$ s.t. $p_S(0)=j$
\item $u_j$ picks a random degree $t$ poly $p_O(x)$ s.t. $p_O(0)=0$
\item $S$ and $u_j$ run a maliciously secure 2-PC where.  I am fairly confident that a generic 2-PC is not needed here, but haven't come up with something better. 
\begin{itemize}
\item $S$ inputs $sk$, $\seckeyserv$, $p_S$, $cp$
\item $u_j$ inputs $p_O$
\item The computation checks that $p_O$ is indeed a poly of degree $t$ with $p_O(0)=0$.  Then, computes $p(x)=p_S(x)+p_O(x)$ by adding the coefficients and then returns to $u_j$ encrypted coefficients of $p(x)$ together with a signature by the server authenticating them.\\
Output (to $u_j$):  \\$c = (Enc_{sk}(\tuple{p_0}{cp_0}), \ldots, Enc_{sk}(\tuple{p_t}{cp_t})), \sigma=Sig_{\seckeyserv}(c)$ 
\end{itemize}
\item $u_j$ includes $(c,\sigma)$ with every copy of the message he sends out.
\end{enumerate}

\paragraph{Complaint:}
To complain about a received message $m$, a party $u_i$ does the following:
\begin{enumerate}
\item $u_i$ first checks that $(c,\sigma)$ included in the message is a valid signature.  If not, he sends the invalid message to $S$.  \anote{May need to include the message in $\sigma$ to prevent replay attack?}
\item $u_i$ uses $Enc_{sk}(\tup{i} ,Enc_{sk}(\tup{i^2}),\ldots,Enc_{sk}(\tup{i^t})$ received from $S$ and $Enc_{sk}(\tuple{p_0}{cp_0}), \ldots, Enc_{sk}(\tuple{p_t}{cp_t})$ received in the message to compute $Enc_{sk}(\tuple{p(i)}{cp(i)})$.
\item $u_i$ sends $Enc_{sk}(\tuple{p(0)}{cp(0)})$ to $S$ who decrypts it and checks that $cp(i)$ is correct.  (Recall that $S$ knows $cp$).
\item If this check passes, $u_i$ does the PIR and write to new location trick to write this complaint to a random location (Note that $S$ will see if he write something other than what he checked in the prior step).  
\end{enumerate}



\bibliography{bib}

\end{document}
